\section*{Chapitre \ref{ch:g9:earth-resources} --- Les ressources de la Terre}

\begin{solution}{exo:g9:earth-resources:1}
(a) \omterm{def:g9:earth-resources:aggregate}{Granulats}. (b) \omterm{def:g9:earth-resources:reserve}{Réserves}. (c) \omterm{def:g9:earth-resources:ore-deposit}{Gisement}. (d) \omterm{def:g9:earth-resources:ore-deposit}{Facteur de concentration}.
\end{solution}

\begin{solution}{exo:g9:earth-resources:2}
(a) Dans les \omterm{def:g3:stones-pebbles-sand:pebble}{graviers} des \omterm{def:g1:landscapes-around-us:water}{rivières}. (b) Dans les fissures, par l'eau
chaude. (c) Dans le \omterm{def:g7:volcanoes:magma}{magma}. (d) C'est la première étape d'un \omterm{def:g9:earth-resources:ore-deposit}{gisement} dans
les \omterm{def:g3:stones-pebbles-sand:pebble}{graviers} d'une \omterm{def:g1:landscapes-around-us:water}{rivière}~: l'or libéré est ensuite emporté par la
\omterm{def:g1:landscapes-around-us:water}{rivière}.
\end{solution}

\begin{solution}{exo:g9:earth-resources:3}
(b), (d), (a), (e), (c).
\end{solution}

\begin{solution}{exo:g9:earth-resources:4}
(a) Non renouvelable~: il met des millions d'années à se former.
(b) Renouvelable~: la \omterm{def:g4:clouds-rain-snow:precipitation}{pluie} recharge ce qui est prélevé. (c) Non
renouvelable~: un \omterm{def:g9:earth-resources:ore-deposit}{gisement} se forme sur des durées très longues. (d) Non
renouvelable à ce rythme~: la nappe baisse et les \omterm{def:g6:water-underground:spring}{puits} finissent à sec.
\end{solution}

\begin{solution}{exo:g9:earth-resources:5}
(a) $4 \div 0.007 \approx 571$, environ $10^3$ (entre $10^2$ et $10^3$,
plus près de $10^3$). (b) $0.5 \div 0.0002 = 2500$, environ $10^3$. (c) $20
\div 5.6 \approx 3.6$, environ 4~: entre $10^0$ et $10^1$.
(d) Le fer est déjà très abondant dans la \omterm{def:g6:inside-the-earth:layers}{croûte}~: un minerai à peine
quelques fois plus riche qu'une \omterm{def:g3:stones-pebbles-sand:rock}{roche} ordinaire suffit.
\end{solution}

\begin{solution}{exo:g9:earth-resources:6}
(a) $19.3 \div 2.5 \approx 7.7$ fois. (b) Or~: $19.3 \times 0.01 =
\qty{0.193}{g}$~; \omterm{def:g3:stones-pebbles-sand:rock}{roche}~: $2.5 \times 0.01 = \qty{0.025}{g}$. (c) Quand
tout le lit se met en mouvement, les grains d'or, lourds, s'enfoncent entre
le \omterm{def:g3:stones-pebbles-sand:sand}{sable} et les \omterm{def:g3:stones-pebbles-sand:pebble}{galets} plus légers, jusqu'à la \omterm{def:g3:stones-pebbles-sand:rock}{roche en place}, où ils
s'arrêtent.
\end{solution}

\begin{solution}{exo:g9:earth-resources:7}
(a) L'or, l'étain et le chrome. (b) Entre $10^1$ et $10^2$, près de
$10^2$. (c) $3000 \div 3.7 \approx 800$ fois. (d) Le chrome (ou l'étain,
ou l'or)~: plus le \omterm{def:g9:earth-resources:ore-deposit}{facteur de concentration} est élevé, plus sont rares les
endroits où la nature a suffisamment concentré le métal.
\end{solution}

\begin{solution}{exo:g9:earth-resources:8}
(a) $18 \times 365 = \qty{6570}{kg}$, environ \qty{6.6}{t}. (b) $5 \times
10^{10}$ tonnes. (c) $25000 \div 18 \approx 1389$ jours, presque quatre
ans. (d) Le \omterm{def:g3:stones-pebbles-sand:sand}{sable} extrait de la \omterm{def:g1:landscapes-around-us:water}{rivière} ne descend plus jusqu'à la \omterm{def:g1:landscapes-around-us:water}{côte}~;
les vagues continuent d'emporter le \omterm{def:g3:stones-pebbles-sand:sand}{sable} de la plage, mais il en arrive
moins pour le remplacer.
\end{solution}

\begin{solution}{exo:g9:earth-resources:9}
(a) Des minuscules plantes et animaux de la \omterm{def:g2:water-on-earth:ocean}{mer}, comme les diatomées,
enfouis dans la \omterm{def:g3:stones-pebbles-sand:sand}{boue} du fond marin. (b) Entre environ 60 et
\qty{120}{\degreeCelsius}. (c) Ils sont plus légers que l'eau, qui remplit
les \omterm{def:g6:water-underground:porosity}{pores} des \omterm{def:g3:stones-pebbles-sand:rock}{roches}. (d) Une couche \omterm{def:g6:water-underground:permeability}{imperméable}, comme une argilite~; ils
s'accumulent sous elle, souvent au sommet d'un pli en voûte, le gaz
au-dessus du pétrole.
\end{solution}

\begin{solution}{exo:g9:earth-resources:10}
(a) $2.9 \times 10^7$ tonnes. (b) $2.9 \times 10^7 \div 100 = 2.9 \times
10^5$ tonnes par an, en gros. (c) Le minerai est 100 fois plus lourd que le
cuivre qu'il contient~: $2.9 \times 10^7 \times 100 = 2.9 \times 10^9$
tonnes. (d) Chaque tonne de minerai ne donne que \qty{10}{kg} de cuivre,
si bien qu'il faut extraire des milliards de tonnes de \omterm{def:g3:stones-pebbles-sand:rock}{roche}.
\end{solution}

\begin{solution}{exo:g9:earth-resources:11}
(a) À 23~millions de tonnes par an, elles auraient duré environ $280 \div 23
\approx 12$ ans, jusqu'au début des années 1980~; or 697~millions de
tonnes ont été extraites, bien plus que ces \omterm{def:g9:earth-resources:reserve}{réserves}. (b) $980 \div 280 = 3.5$
fois. (c) $980 \div
23 \approx 43$ ans. (d) Par exemple~: on découvre de nouveaux \omterm{def:g9:earth-resources:ore-deposit}{gisements}~;
de meilleures techniques ou des prix plus élevés rendent rentables des
\omterm{def:g9:earth-resources:ore-deposit}{gisements} plus pauvres.
\end{solution}

\begin{solution}{exo:g9:earth-resources:12}
(a) $5 \times 10^1$ ans. (b) $2 \times 10^3 \div 50 = 40$ fois. (c) $3
\times 10^8 \div (5 \times 10^1) = 6 \times 10^6$ fois. (d) Le charbon a
mis des millions de fois plus de temps à se former que nous n'en mettons à
l'utiliser~: la nature ne le renouvelle pas à notre échelle de temps.
\end{solution}

\begin{solution}{pb:g9:earth-resources:1}
\textbf{1.} $100 \times 1000 = \num{100000}$~kg $= 1 \times 10^5$~kg.

\textbf{2.} $2000 \times 300 = \num{600000}$~t $= 6 \times 10^5$~t.

\textbf{3.} $6 \times 10^5 \div 10 = 6 \times 10^4$~t par an.

\textbf{4.} $0.40 \times 6 \times 10^4 = 2.4 \times 10^4$~t de \omterm{def:g3:stones-pebbles-sand:sand}{sable} par
an.

\textbf{5.} $1200 \times 500 \times 5 = 3 \times 10^6$~m$^3$.

\textbf{6.} Les deux tiers~: $2 \times 10^6$~m$^3$.

\textbf{7.} $2 \times 10^6 \times 1.6 = 3.2 \times 10^6$~t.

\textbf{8.} $\num{24000} + \num{136000} = \num{160000}$~t $= 1.6 \times
10^5$~t par an.

\textbf{9.} $\num{160000} \div 2000 = 80$ fois plus.

\textbf{10.} Non~: la carrière se vide 80 fois plus vite que la \omterm{def:g1:landscapes-around-us:water}{rivière} ne
renouvelle le \omterm{def:g3:stones-pebbles-sand:sand}{sable}. À l'échelle de la ville, ce \omterm{def:g3:stones-pebbles-sand:sand}{sable} est une ressource
épuisable.

\textbf{11.} $\num{24000} \div 4 = \num{6000}$~t de \omterm{def:g3:stones-pebbles-sand:sand}{sable} par an.

\textbf{12.} $3.2 \times 10^6 \div (1.6 \times 10^5) = 20$~: \textbf{20 ans}.
\end{solution}
